\section{Desenvolvimento}

\begin{frame}{Metodo}
  Descreva o percurso com o detalhe suficiente para que alguem possa repeti-lo,
  apoiando-se na literatura quando couber \cite{exemplo2024}.

  \vspace{1em}
  \begin{enumerate}
    \item Coleta dos dados.
    \item Tratamento e analise.
    \item Validacao dos resultados.
  \end{enumerate}
\end{frame}

\begin{frame}{Resultados}
  \begin{table}
    \centering
    \begin{tabular}{lrr}
      \toprule
      \textbf{Condicao} & \textbf{Media} & \textbf{Desvio} \\
      \midrule
      Controle      & 12,4 & 1,8 \\
      Experimental  & 18,7 & 2,1 \\
      \bottomrule
    \end{tabular}
  \end{table}
  \fonte{Elaborado pelo autor (2026).}
\end{frame}

\begin{frame}{Figuras}
  % Coloque a imagem em figures/ e descomente:
  %
  % \begin{figure}
  %   \centering
  %   \includegraphics[width=0.7\textwidth]{exemplo.png}
  %   \caption{Descricao breve do que a figura mostra}
  % \end{figure}
  % \fonte{Elaborado pelo autor (2026).}

  Uma imagem que ocupa o slide inteiro comunica melhor do que quatro graficos
  espremidos. Prefira \alert{um grafico por slide}.
\end{frame}
